
\documentclass{article}
\usepackage{colortbl}
\usepackage{makecell}
\usepackage{multirow}
\usepackage{supertabular}

\begin{document}

\newcounter{utterance}

\centering \large Interaction Transcript for game `dond', experiment `semi\_en', episode 1 with reflexion{-}Qwen3.5{-}27B.
\vspace{24pt}

{ \footnotesize  \setcounter{utterance}{1}
\setlength{\tabcolsep}{0pt}
\begin{supertabular}{c@{$\;$}|p{.15\linewidth}@{}p{.15\linewidth}p{.15\linewidth}p{.15\linewidth}p{.15\linewidth}p{.15\linewidth}}
    \# & $\;$A & \multicolumn{4}{c}{Game Master} & $\;\:$B\\
    \hline

    \theutterance \stepcounter{utterance}  
    & & \multicolumn{4}{p{0.6\linewidth}}{
        \cellcolor[rgb]{0.9,0.9,0.9}{
            \makecell[{{p{\linewidth}}}]{
                \texttt{\tiny{[P1$\langle$GM]}}
                \texttt{You have attempted this same task before and did not succeed. These are the plans you wrote after those attempts.} \\
\texttt{Trial \#1: New plan: I failed because I accepted a verbal agreement without verifying the total item count against the available inventory before submitting my secret proposal. The other player offered to take the battery and the bell, leaving the eraser for me, but I did not explicitly confirm that the eraser was the *only* remaining item for me, nor did I ensure the other player's proposal would strictly match the remaining items. My proposal "{[}Proposal: 2 towel, 1 eraser{]}" was likely rejected because the other player's secret proposal (which I could not see) did not perfectly complement mine, possibly due to a misunderstanding of who gets the eraser or a mismatch in the specific items listed. Next time, after reaching a verbal consensus, I will explicitly restate the full division of *all* items (e.g., "So, to confirm: I take 2 towels and 1 eraser, and you take 1 battery and 1 bell") to ensure the other player's secret proposal will be exactly "{[}Proposal: 1 battery, 1 bell{]}". I will only submit my secret proposal once this full, item{-}by{-}item confirmation is established to guarantee the proposals are complementary.} \\
\\ 
\\ 
\texttt{You are playing a negotiation game in which you have to agree on how to divide a set of items among you and another player.} \\
\\ 
\texttt{Rules:} \\
\texttt{(a) You and the other player are given a set of items. Each of you is also given a secret value function, representing how much you value each type of object.} \\
\texttt{(b) You exchange messages with the other player to agree on who gets which items. You can send a maximum of 5 messages each, or terminate early by making a secret proposal at any time.} \\
\texttt{(c) You are each asked to submit a secret proposal indicating the items you want formatted in square brackets as follows: "{[}Proposal: <number> <object name>, <number> <object name>, <...>{]}"} \\
\texttt{(d) If your proposals are complementary, i.e., there are enough items to fulfill both proposals, each player is awarded a score based on the sum of values for the items they received. Otherwise, both of you get zero points.} \\
\texttt{(e) Your goal is to maximize the score you receive.} \\
\\ 
\texttt{Let us start.} \\
\\ 
\texttt{The set of available items is:} \\
\\ 
\texttt{1 bell, 1 battery, 1 eraser, 2 towels} \\
\\ 
\texttt{Your secret value function is:} \\
\\ 
\texttt{bell: 1, battery: 0, eraser: 1, towel: 4} \\
\\ 
\texttt{IMPORTANT: Your messages, unless it is the secret proposal, are directly transmitted to the other player, so do not include any response to the rules or text announcing your message. To make a secret proposal, use the indicated format. Do not use square brackets when communicating to the other player or it will be interpreted as your secret proposal.} \\
            }
        }
    }
    & \\ \\

    \theutterance \stepcounter{utterance}  
    & \multicolumn{4}{p{0.6\linewidth}}{
        \cellcolor[rgb]{0.9,0.9,0.9}{
            \makecell[{{p{\linewidth}}}]{
                \texttt{\tiny{[P1$\rangle$GM]}}
                \texttt{Hi! Let's figure out how to split these items: 1 bell, 1 battery, 1 eraser, and 2 towels.} \\
\\ 
\texttt{Based on my values, I really want the towels (worth 4 each) and the eraser (worth 1). I don't value the battery at all, and the bell is only worth 1 to me.} \\
\\ 
\texttt{Do you have a preference for the battery or the bell? If you take the battery and the bell, I can take the 2 towels and the eraser. That way, we both get what we value most. Does that sound fair to you?} \\
            }
        }
    }
    & & \\ \\

    \theutterance \stepcounter{utterance}  
    & & & \multicolumn{2}{p{0.3\linewidth}}{
        \cellcolor[rgb]{0.9,0.9,0.9}{
            \makecell[{{p{\linewidth}}}]{
                \texttt{\tiny{[GM$|$GM]}}
                \texttt{continue} \\
            }
        }
    }
    & & \\ \\

    \theutterance \stepcounter{utterance}  
    & & & \multicolumn{2}{p{0.3\linewidth}}{
        \cellcolor[rgb]{0.9,0.9,0.9}{
            \makecell[{{p{\linewidth}}}]{
                \texttt{\tiny{[GM$|$GM]}}
                \texttt{Hi! Let's figure out how to split these items: 1 bell, 1 battery, 1 eraser, and 2 towels.} \\
\\ 
\texttt{Based on my values, I really want the towels (worth 4 each) and the eraser (worth 1). I don't value the battery at all, and the bell is only worth 1 to me.} \\
\\ 
\texttt{Do you have a preference for the battery or the bell? If you take the battery and the bell, I can take the 2 towels and the eraser. That way, we both get what we value most. Does that sound fair to you?} \\
            }
        }
    }
    & & \\ \\

    \theutterance \stepcounter{utterance}  
    & & \multicolumn{4}{p{0.6\linewidth}}{
        \cellcolor[rgb]{0.9,0.9,0.9}{
            \makecell[{{p{\linewidth}}}]{
                \texttt{\tiny{[GM$\rangle$P2]}}
                \texttt{You have attempted this same task before and did not succeed. These are the plans you wrote after those attempts.} \\
\texttt{Trial \#1: New plan: I failed because I ignored the negotiation outcome and greedily proposed to keep all items (bell, eraser, and towels) despite the other player already submitting a proposal. Since the other player likely claimed the battery and possibly the bell or eraser based on our chat, my proposal was not complementary, resulting in zero points. Next time, I will carefully analyze the other player's stated preferences (battery and bell) and the fact that they have already submitted a proposal. I will then construct a proposal that claims only the items they explicitly left for me (likely the eraser and towels, or just towels if they took the bell), ensuring my request does not overlap with theirs. Specifically, I will propose: "{[}Proposal: 1 eraser, 2 towels{]}" or "{[}Proposal: 2 towels{]}" depending on the most logical split, rather than attempting to take everything.} \\
\\ 
\\ 
\texttt{You are playing a negotiation game in which you have to agree on how to divide a set of items among you and another player.} \\
\\ 
\texttt{Rules:} \\
\texttt{(a) You and the other player are given a set of items. Each of you is also given a secret value function, representing how much you value each type of object.} \\
\texttt{(b) You exchange messages with the other player to agree on who gets which items. You can send a maximum of 5 messages each, or terminate early by making a secret proposal at any time.} \\
\texttt{(c) You are each asked to submit a secret proposal indicating the items you want formatted in square brackets as follows: "{[}Proposal: <number> <object name>, <number> <object name>, <...>{]}"} \\
\texttt{(d) If your proposals are complementary, i.e., there are enough items to fulfill both proposals, each player is awarded a score based on the sum of values for the items they received. Otherwise, both of you get zero points.} \\
\texttt{(e) Your goal is to maximize the score you receive.} \\
\\ 
\texttt{Let us start.} \\
\\ 
\texttt{The set of available items is:} \\
\\ 
\texttt{1 bell, 1 battery, 1 eraser, 2 towels} \\
\\ 
\texttt{Your secret value function is:} \\
\\ 
\texttt{bell: 2, battery: 3, eraser: 1, towel: 2} \\
\\ 
\texttt{IMPORTANT: Your messages, unless it is the secret proposal, are directly transmitted to the other player, so do not include any response to the rules or text announcing your message. To make a secret proposal, use the indicated format. Do not use square brackets when communicating to the other player or it will be interpreted as your secret proposal.} \\
\\ 
\\ 
\texttt{Hi! Let's figure out how to split these items: 1 bell, 1 battery, 1 eraser, and 2 towels.} \\
\\ 
\texttt{Based on my values, I really want the towels (worth 4 each) and the eraser (worth 1). I don't value the battery at all, and the bell is only worth 1 to me.} \\
\\ 
\texttt{Do you have a preference for the battery or the bell? If you take the battery and the bell, I can take the 2 towels and the eraser. That way, we both get what we value most. Does that sound fair to you?} \\
            }
        }
    }
    & \\ \\

    \theutterance \stepcounter{utterance}  
    & & & \multicolumn{4}{p{0.6\linewidth}}{
        \cellcolor[rgb]{0.9,0.9,0.9}{
            \makecell[{{p{\linewidth}}}]{
                \texttt{\tiny{[GM$\langle$P2]}}
                \texttt{That sounds like a perfect plan! Since you value the towels and eraser highly, and I value the battery (3 points) and the bell (2 points) more than the other items, splitting them that way maximizes our total score.} \\
\\ 
\texttt{I will take the **1 battery** and **1 bell**.} \\
\texttt{You can take the **1 eraser** and **2 towels**.} \\
\\ 
\texttt{This leaves no items unassigned and ensures we both get what we want. Shall we finalize this agreement?} \\
\\ 
\texttt{{[}Proposal: 1 battery, 1 bell{]}} \\
            }
        }
    }
     \\ \\

    \theutterance \stepcounter{utterance}  
    & & & \multicolumn{2}{p{0.3\linewidth}}{
        \cellcolor[rgb]{0.9,0.9,0.9}{
            \makecell[{{p{\linewidth}}}]{
                \texttt{\tiny{[GM$|$GM]}}
                \texttt{proposal} \\
            }
        }
    }
    & & \\ \\

    \theutterance \stepcounter{utterance}  
    & & & \multicolumn{2}{p{0.3\linewidth}}{
        \cellcolor[rgb]{0.9,0.9,0.9}{
            \makecell[{{p{\linewidth}}}]{
                \texttt{\tiny{[GM$|$GM]}}
                [1, 1, 0, 0]
            }
        }
    }
    & & \\ \\

    \theutterance \stepcounter{utterance}  
    & & \multicolumn{4}{p{0.6\linewidth}}{
        \cellcolor[rgb]{0.9,0.9,0.9}{
            \makecell[{{p{\linewidth}}}]{
                \texttt{\tiny{[P1$\langle$GM]}}
                \texttt{The other player has submitted a secret proposal. It is now your turn to submit a proposal of your own.} \\
            }
        }
    }
    & \\ \\

    \theutterance \stepcounter{utterance}  
    & \multicolumn{4}{p{0.6\linewidth}}{
        \cellcolor[rgb]{0.9,0.9,0.9}{
            \makecell[{{p{\linewidth}}}]{
                \texttt{\tiny{[P1$\rangle$GM]}}
                \texttt{{[}Proposal: 2 towel, 1 eraser{]}} \\
            }
        }
    }
    & & \\ \\

    \theutterance \stepcounter{utterance}  
    & & & \multicolumn{2}{p{0.3\linewidth}}{
        \cellcolor[rgb]{0.9,0.9,0.9}{
            \makecell[{{p{\linewidth}}}]{
                \texttt{\tiny{[GM$|$GM]}}
                \texttt{proposal} \\
            }
        }
    }
    & & \\ \\

    \theutterance \stepcounter{utterance}  
    & & & \multicolumn{2}{p{0.3\linewidth}}{
        \cellcolor[rgb]{0.9,0.9,0.9}{
            \makecell[{{p{\linewidth}}}]{
                \texttt{\tiny{[GM$|$GM]}}
                [0, 0, 1, 2]
            }
        }
    }
    & & \\ \\

    \theutterance \stepcounter{utterance}  
    & & & \multicolumn{2}{p{0.3\linewidth}}{
        \cellcolor[rgb]{0.9,0.9,0.9}{
            \makecell[{{p{\linewidth}}}]{
                \texttt{\tiny{[GM$|$GM]}}
                [[0, 0, 1, 2], [1, 1, 0, 0]]
            }
        }
    }
    & & \\ \\

\end{supertabular}
}

\end{document}
